\subsection{余均衡子的定义}

设 H 为一个箭图，G 为一个群胚，\(\varphi\) 、\(\psi\) 为从 H 到 G 的箭图态射。以 \(\mathrm{Coh}(\varphi,\psi)\) 表示偶 \((\varphi,\psi)\) 的上同伦子（II, 第 185 页, 定义 3），以 \(\alpha\colon\mathrm{G}\to\mathrm{Coh}(\varphi,\psi)\) 表示典范群胚态射，以 \(h\) 表示联系 \(\alpha\circ\varphi\) 与 \(\alpha\circ\psi\) 的典范同伦。

设 \(\operatorname{Coeg}(\varphi,\psi)\) 为由 \(\operatorname{Coh}(\varphi,\psi)\) 通过缩并属于 h 的像的箭而得到的群胚（II, 第 196 页, 第 1 号）；以 \(\beta\colon\operatorname{Coh}(\varphi,\psi)\to\operatorname{Coeg}(\varphi,\psi)\) 表示典范态射，并令 \(\gamma=\beta\circ\alpha\) 。

定义 2. —— 我们称群胚 \(\operatorname{Coeg}(\varphi,\psi)\) 为偶 \((\varphi,\psi)\) 的余均衡子；群胚态射 \(\gamma\) 称为从 G 到 \(\operatorname{Coeg}(\varphi,\psi)\) 的典范态射。

命题 3. —— 偶 (Coeg(φ,ψ),γ) 具有下列泛性质：

\begin{enumerate}
\def\labelenumi{\alph{enumi})}
\item
  我们有 \(\gamma \circ \varphi = \gamma \circ \psi\) 。
\item
  设 G′ 为一个群胚，\(\theta\colon G\to G'\) 为满足 \(\theta\circ\varphi=\theta\circ\psi\) 的一个群胚态射。存在唯一的群胚态射 \(\overline{\theta}\colon\operatorname{Coeg}(\varphi,\psi)\to G'\) ，满足 \(\overline{\theta}\circ\gamma=\theta\) 。
\end{enumerate}

设 \(a\) 为 H 的一个顶点。箭 \(h(a)\) （\(\mathrm{Coh}(\varphi, \psi)\) 的）连接 \(\alpha(\varphi(a))\) 与 \(\alpha(\psi(a))\) 。根据群胚 \(\mathrm{Coeg}(\varphi, \psi)\) 的定义，箭 \(\beta(h(a))\) 的起点与靶相等；因此我们有 \(\gamma(\varphi(a)) = \gamma(\psi(a))\) 。

设 f 为 H 的一条箭；若以 a 表示其起点，以 b 表示其靶，则我们有

\[
\alpha (\varphi (f)) \cdot h (b) = h (a) \cdot \alpha (\psi (b)).
\]

对此等式的两边取在 \(\beta\) 下的像，我们得到关系式 \(\gamma(\varphi(f)) = \gamma(\psi(f))\) 。这就证明了断言 \(a\) )。

我们来证明 \(b\) )。映射 \(\eta\colon\mathrm{Som}(\mathrm{H})\to\mathrm{Fl}(\mathrm{G}')\) 把 H 的每个顶点 \(a\) 对应到 G′ 的箭 \(e_{\theta(\varphi(a))}\) ，它是连接 \(\theta\circ\varphi\) 与 \(\theta\circ\psi\) 的一个同伦。由上同伦子的泛性质（II, 第 185 页, 命题 3），存在唯一的群胚态射 \(\theta_1\colon\mathrm{Coh}(\varphi,\psi)\to\mathrm{G}'\) ，满足 \(\theta_1\circ\alpha=\theta\) 以及 \(\theta_1(h(a))=e_{\theta(\varphi(a))}\) （对 H 的每个顶点 \(a\) ）。由 II, 第 176 页的命题 7，后一性质蕴含存在唯一的群胚态射 \(\overline{\theta}\colon\mathrm{Coeg}(\varphi,\psi)\to\mathrm{G}'\) ，满足 \(\overline{\theta}\circ\beta=\theta_1\)。于是有 \(\overline{\theta}\circ\gamma=\theta_1\circ\alpha=\theta\)。

反之，若 \(\overline{\theta}'\colon \mathrm{Coeg}(\varphi,\psi)\to \mathrm{G}'\) 是满足 \(\overline{\theta}'\circ\gamma=\theta\) 的一个群胚态射，则我们有 \((\overline{\theta}'\circ\beta)\circ\alpha=(\overline{\theta}\circ\beta)\circ\alpha\) ，从而由 II, 第 185 页, 命题 3 得 \(\overline{\theta}'\circ\beta=\overline{\theta}\circ\beta\) ，再由 II, 第 176 页的命题 7 得 \(\overline{\theta}'=\overline{\theta}\) 。这就证明了 \(\overline{\theta}\) 的唯一性，从而断言 \(b\) ) 得证。

推论. —— 群胚 Coeg(φ,ψ) 由态射 γ 的像生成。

设 C 为由 G 的像生成的 Coeg( \(\varphi,\psi\) ) 的子群胚；以 i 表示从 C 到 Coeg( \(\varphi,\psi\) ) 的典范态射，以 \(\theta\colon G\to C\) 表示满足 \(i\circ\theta=\gamma\) 的态射。由命题 2，存在唯一的群胚态射 \(\overline{\theta}\colon\) Coeg( \(\varphi,\psi\) ) \(\to\) C ，满足 \(\overline{\theta}\circ\gamma=\theta\) 。于是，\(\gamma=i\circ\overline{\theta}\circ\gamma\) ，因此 \(i\circ\overline{\theta}\) 是 Coeg( \(\varphi,\psi\) ) 的恒同态射（同前）。特别地，映射 Som(i) 与 Fl(i) 都是满射的，从而 C = Coeg( \(\varphi,\psi\) )。

注记. —— 1) Coeg( \(\varphi,\psi\) ) 的顶点集是箭图

\[
\Gamma = (\operatorname{Som} (\mathrm{G}), \operatorname{Som} (\mathrm{H}), \operatorname{Som} (\varphi), \operatorname{Som} (\psi))
\]

（II, 第 197 页）的连通分支所成的集。换言之，它是 \(\operatorname{Som}(\mathbf{G})\) 关于使 \(\varphi(a)\) 与 \(\psi(a)\) 等价（对所有 \(a \in \operatorname{Som}(\mathbf{G})\) ）的最细等价关系的商集。

\begin{enumerate}
\def\labelenumi{\arabic{enumi})}
\setcounter{enumi}{1}
\tightlist
\item
  映射 \(\mathrm{Orb}(\gamma)\) 由 \(\gamma\) 通过取轨道导出，它通过取商定义了一个从偶 \((\varphi,\psi)\) 的框架的连通分支集到余均衡子的轨道集的双射。这由 II, 第 185 页, 命题 4 以及由 \(\beta\) 通过取轨道导出的映射是双射的这一事实（II, 第 170 页, 注 1）得出。
\end{enumerate}

因此，从 Som(G) 到 \(\mathrm{Orb}(\mathrm{Coeg}(\varphi,\psi))\) 的、由 \(\gamma\) 导出的映射把 \(\mathrm{Coeg}(\varphi,\psi)\) 的轨道集与 Som(G) 关于由偶 \((\varphi(x),\psi(x))\) （对 \(x\in\mathrm{Som(H)}\) ）以及偶 \((o(f),t(f))\) （对 \(f\in\mathrm{Fl(G)}\) ）生成的等价关系的商集等同起来。

\begin{enumerate}
\def\labelenumi{\arabic{enumi})}
\setcounter{enumi}{2}
\tightlist
\item
  从 Fl(G) 到 Fl(Coeg(φ,ψ)) 的映射 Fl(γ) 一般不是满射的。
\end{enumerate}

